% =========================================================
% Section: Carathéodory Factorization and Chain Rule
% =========================================================

\begin{tcolorbox}[colback=gray!8, colframe=gray!40, title=\textbf{Carathéodory and Chain Rule --- toolkit}]
\textbf{Concept family:} The Carathéodory factorization of the derivative, its characterization theorem, and the Chain Rule.

\medskip
\textbf{Atomic items in this section:}
\begin{itemize}
  \item Definition: Carathéodory Factorization
  \item Theorem: Carathéodory Characterization of Differentiability
  \item Theorem: Chain Rule
\end{itemize}
\end{tcolorbox}

% ---------------------------------------------------------

\begin{tcolorbox}[colback=defbox, colframe=defborder]
\begin{definition}[Carathéodory Factorization]
\label{def:caratheodory-factorization}
Let $I \subseteq \mathbb{R}$ be an interval, let $f : I \to \mathbb{R}$, and let $c \in I$. We say that $f$ admits a \emph{Carathéodory factorization} at $c$ if there exists a function $\phi : I \to \mathbb{R}$, continuous at $c$, such that
\[
f(x) = f(c) + (x-c)\,\phi(x) \qquad \text{for all } x \in I.
\]
In this case, $f$ is differentiable at $c$ and $f'(c) = \phi(c)$.
\end{definition}
\end{tcolorbox}

\begin{remark*}[Standard quantified statement]
\[
\exists \phi : I \to \mathbb{R}\;
\bigl(\operatorname{ContinuousAt}(\phi,c)
\;\land\;
\forall x \in I\;[f(x) = f(c)+(x-c)\phi(x)]\bigr).
\]
\end{remark*}

\begin{remark*}[Definition predicate reading]
The existence of a Carathéodory factorization at $c$ is equivalent to $\operatorname{DifferentiableAt}(f,c)$ by the characterization theorem below. The factorizing function is explicit:
\[
\phi(x) =
\begin{cases}
\dfrac{f(x)-f(c)}{x-c} & x \neq c, \\[6pt]
f'(c) & x = c.
\end{cases}
\]
Differentiability of $f$ at $c$ is precisely the statement that this piecewise-defined $\phi$ is continuous at $c$.
\end{remark*}

\begin{remark*}[Interpretation]
The Carathéodory factorization replaces the difference quotient --- a function with a removable discontinuity at $c$ --- with a function $\phi$ that is genuinely continuous at $c$. The factorization $f(x)-f(c) = (x-c)\phi(x)$ holds everywhere on $I$, including at $x=c$ where both sides equal zero. The derivative is recovered as $\phi(c)$ without ever dividing by zero. This reframing converts analytical limit conditions into topological continuity conditions, which compose and factor more cleanly. It is the preferred tool for proving the product rule, quotient rule, and chain rule, because continuity of a product or composition is easier to establish than the existence of a limit of a quotient.

\medskip
\textbf{Carathéodory vs.\ Lipschitz.} Both control the increment $f(x)-f(c)$ by the factor $(x-c)$, but differently. The Carathéodory factorization requires the slope function $\phi(x)$ to extend \emph{continuously} to $c$; it controls the \emph{behaviour} of the slope, not just its size. The Lipschitz condition $|f(x)-f(c)| \leq L|x-c|$ only bounds the size of the increment. Carathéodory implies Lipschitz: if $\phi$ is continuous at $c$ it is bounded near $c$, so $|f(x)-f(c)| = |\phi(x)||x-c| \leq K|x-c|$ for some $K$ near $c$. The converse fails: $f(x)=|x|$ is Lipschitz on $\mathbb{R}$ with constant $1$ but the slope function takes values $\pm 1$ on either side of $0$ with no continuous extension.
\end{remark*}

\begin{remark*}[Dependencies]
\begin{itemize}
  \item \hyperref[def:derivative-at-a-point]{Definition: Derivative at a Point}
  \item \hyperref[def:continuous-at]{Definition: Continuous at a Point}
\end{itemize}
\end{remark*}

% ---------------------------------------------------------

\begin{theorem}[Carathéodory Characterization of Differentiability]
\label{thm:caratheodory-characterization-of-differentiability}
Let $I \subseteq \mathbb{R}$ be an interval, let $f : I \to \mathbb{R}$, and let $c \in I$. The following are equivalent:
\begin{enumerate}[label=(\roman*)]
\item $f$ is differentiable at $c$.
\item There exists a function $\phi : I \to \mathbb{R}$, continuous at $c$, such that
\[
f(x) = f(c) + (x-c)\,\phi(x) \qquad \text{for all } x \in I.
\]
\end{enumerate}
In this case, $\phi(c) = f'(c)$.
\hyperref[prf:caratheodory-theorem]{Go to proof.}
\end{theorem}

\begin{remark*}[Standard quantified statement]
\begin{align*}
\operatorname{DifferentiableAt}(f,c)
\;&\iff\;
\exists\phi:I\to\mathbb{R}\;
\bigl(\operatorname{ContinuousAt}(\phi,c)\\
&\quad\land\;\forall x \in I\;[f(x)=f(c)+(x-c)\phi(x)]\bigr).
\end{align*}
In this case $\phi(c) = f'(c)$.
\end{remark*}

\begin{remark*}[Interpretation]
This theorem establishes that differentiability and the existence of a Carathéodory factorization are exactly the same condition. The proof is a two-direction construction: in the forward direction, define $\phi$ as the difference quotient extended by $f'(c)$ at $c$ and verify continuity from the definition of the derivative; in the backward direction, recover the difference quotient from $\phi$ and read off its limit as $\phi(c)$. The theorem's power is that it converts every future differentiability argument into a continuity argument: to show $g \circ f$ is differentiable, find a continuous Carathéodory function for it; composing the Carathéodory functions of $f$ and $g$ yields one automatically.
\end{remark*}

\begin{remark*}[Dependencies]
\begin{itemize}
  \item \hyperref[def:caratheodory-factorization]{Definition: Carathéodory Factorization}
  \item \hyperref[def:derivative-at-a-point]{Definition: Derivative at a Point}
  \item \hyperref[def:continuous-at]{Definition: Continuous at a Point}
\end{itemize}
\end{remark*}

% ---------------------------------------------------------

\begin{theorem}[Chain Rule]
\label{thm:chain-rule}
Let $I, J \subseteq \mathbb{R}$ be intervals, let $f : I \to \mathbb{R}$ and $g : J \to \mathbb{R}$, and suppose that $f(I) \subseteq J$. Let $c \in I$. If $f$ is differentiable at $c$ and $g$ is differentiable at $f(c)$, then the composition $g \circ f$ is differentiable at $c$, and
\[
(g \circ f)'(c) = g'(f(c))\,f'(c).
\]
\hyperref[prf:chain-rule]{Go to proof.}
\end{theorem}

\begin{remark*}[Standard quantified statement]
\begin{align*}
&\operatorname{DifferentiableAt}(f,c) \;\land\; \operatorname{DifferentiableAt}(g,f(c))\\
&\quad\Rightarrow\;
\operatorname{DerivativeAt}(g\circ f,\,c,\,g'(f(c))\cdot f'(c)).
\end{align*}
\end{remark*}

\begin{remark*}[Interpretation]
The chain rule computes the derivative of a composition by multiplying the derivatives at the respective points: the outer function $g$ is differentiated at the intermediate value $f(c)$, and the inner function $f$ is differentiated at $c$. The Carathéodory proof is the cleanest: let $\phi_f$ and $\phi_g$ be the Carathéodory functions of $f$ at $c$ and $g$ at $f(c)$ respectively. Then
\[
g(f(x)) - g(f(c)) = \phi_g(f(x))\cdot(f(x)-f(c)) = \phi_g(f(x))\cdot\phi_f(x)\cdot(x-c).
\]
The function $x \mapsto \phi_g(f(x))\cdot\phi_f(x)$ is continuous at $c$ --- the product of two continuous functions --- and its value at $c$ is $g'(f(c))\cdot f'(c)$. By the characterization theorem, $g\circ f$ is differentiable with the stated derivative. The naive proof via difference quotients fails because $f(x) = f(c)$ for $x$ near $c$ is possible, causing division by zero; Carathéodory avoids this entirely.
\end{remark*}

\begin{remark*}[Dependencies]
\begin{itemize}
  \item \hyperref[thm:caratheodory-characterization-of-differentiability]{Theorem: Carathéodory Characterization of Differentiability}
  \item \hyperref[thm:differentiable-implies-continuous]{Theorem: Differentiable Implies Continuous}
  \item \hyperref[thm:continuity-composition]{Theorem: Continuity of Composite Functions}
  \item \hyperref[thm:continuity-product]{Theorem: Product of Continuous Functions is Continuous}
\end{itemize}
\end{remark*}
